\chapter*{ABSTRAK}
\addcontentsline{toc}{chapter}{ABSTRAK}

\begin{center}
    \textbf{IMPLEMENTASI PURWARUPA SISTEM ASESMEN ADAPTIF KOLABORATIF BERBASIS \textit{ELO RATING} DENGAN MITIGASI \textit{OVERPERSONALIZATION}} \\
    \vspace{0.4cm}
    oleh \\
    \vspace{0.4cm}
    Dama Dhananjaya Daliman\\
    18222047\\
\end{center}

\begin{spacing}{1.0}
Perkembangan sistem asesmen adaptif bertujuan untuk mengoptimalkan jalur \linebreak belajar individu melalui personalisasi konten. Namun, ketergantungan pada algoritma personalisasi yang agresif berisiko memicu fenomena \textit{overpersonalization}, yang menciptakan \textit{filter bubbles} akademik dan mengisolasi pelajar dari keragaman perspektif. Kondisi ini dapat menyebabkan pengambilan sampel informasi yang selektif serta stagnasi pembelajaran (\textit{learning plateau}).

Penelitian ini bertujuan untuk merancang dan mengimplementasikan purwarupa sistem asesmen \linebreak adaptif kolaboratif. Sistem diimplementasikan menggunakan arsitektur \textit{client-server} dengan algoritma \textit{Modified Elo Rating} untuk memodelkan kemampuan kognitif \linebreak pelajar. Selanjutnya, mekanisme mitigasi diimplementasikan melalui deteksi stagnasi berbasis \linebreak variansi perubahan skor ($\Delta\theta$) dan algoritma \textit{Constraint-Based Re-ranking} untuk \linebreak memasangkan pelajar dengan sejawat yang memiliki tingkat kompetensi heterogen (\textit{Cohen's d} $\geq 0{,}5$). Terakhir, evaluasi sistem dilakukan melalui \textit{Controlled Lab Study} yang melibatkan sepuluh partisipan yang dibagi ke dalam kelompok intervensi dan kontrol pada domain pembelajaran SQL.

Hasil pengujian menunjukkan bahwa kelompok intervensi mencapai peningkatan hasil belajar yang signifikan dengan nilai \textit{Normalized Learning Gain} (NLG) antara 0,57 hingga 0,69 (kategori sedang ke tinggi), sementara kelompok kontrol \linebreak mengalami penurunan performa (NLG -0,21). Analisis trajektori kemahiran \linebreak mengindikasikan intervensi kolaboratif mampu mengembalikan pertumbuhan belajar dengan rata-rata perbedaan gradien sebesar +0,67 setelah titik stagnasi terdeteksi. Algoritma \textit{re-ranking} juga terbukti efektif menciptakan pasangan heterogen dengan rata-rata \textit{Cohen's d} sebesar 1,186. Dari penelitian ini, disimpulkan bahwa integrasi \linebreak intervensi sosial yang proaktif dan terautomasi menunjukkan indikasi efektif untuk memitigasi \textit{overpersonalization} dalam lingkungan pembelajaran adaptif. Ke depannya penelitian dapat dikembangkan dengan rekalibrasi parameter stagnasi dan \linebreak integrasi LLM untuk \textit{scaffolding peer review}. 

\vspace{0.4cm}
\noindent \textbf{Kata kunci:} Asesmen Adaptif, Pembelajaran Kolaboratif, \textit{Elo rating}, \textit{overpersonalization}\par
\end{spacing}

\newpage

% \chapter*{ABSTRACT}
% \addcontentsline{toc}{chapter}{ABSTRACT}

% \begin{center}
%     \textbf{IMPLEMENTATION OF A COLLABORATIVE ADAPTIVE ASSESSMENT SYSTEM PROTOTYPE BASED ON ELO RATING WITH OVERPERSONALIZATION MITIGATION} \\
%     \vspace{0.4cm}
%     by \\
%     \vspace{0.4cm}
%     Dama Dhananjaya Daliman\\
%     18222047\\
% \end{center}

% \begin{spacing}{1.0}
% The development of adaptive assessment systems aims to optimize individual learning paths through content personalization. However, reliance on aggressive personalization algorithms risks triggering the overpersonalization phenomenon, which creates academic filter bubbles and isolates learners from diverse perspectives. This condition can lead to selective information sampling and cognitive stagnation (learning plateau). This research aims to design and implement a collaborative adaptive assessment system prototype named Equilibria, which integrates social features as a mitigation effort against overpersonalization effects.

% The system is implemented using a client-server architecture with a Modified Elo Rating algorithm for cognitive ability tracking. The mitigation mechanism is realized through stagnation detection based on the variance of score changes ($\Delta\theta$) and a Constraint-Based Re-ranking algorithm to pair students with peers of heterogeneous competency levels ($d \geq 0.5$). Evaluation was conducted through a Controlled Lab Study involving 10 participants in intervention and control groups within the SQL learning domain.

% The test results show that the intervention group achieved significant learning gains with Normalized Learning Gain (NLG) values between 0.57 and 0.69 (medium to high category), while the control group experienced a performance decline (NLG -0.21). Proficiency trajectory analysis proves that collaborative intervention successfully restored learning growth momentum with an average positive gradient difference of +0.67 after the stagnation point was detected. The re-ranking algorithm also proved effective in creating heterogeneous pairs with an average Cohen's d of 1.186. This study concludes that the integration of proactive and automated social intervention is an effective strategy to break academic isolation in adaptive learning environments.

% \vspace{0.4cm}
% \noindent \textbf{Keywords:} Adaptive Assessment, Collaborative Learning, Elo rating, overpersonalization, Filter Bubble.
% \end{spacing}
